\begin{frame}{여러 시드: 한 곡선으로 판단하지 마라}
  \begin{itemize}
    \item 위 숫자(21.15 $\to$ 103.85)는 \textbf{시드 0} 한 시드의 대표값.
    \item 신경망 RL은 \textbf{세 곳의 무작위성}을 동시에 안고 있음:
      초기화, $\varepsilon$-탐색의 행동 선택, 버퍼 샘플링
      $\Rightarrow$ \textbf{시드마다 곡선 모양이 다름} (500 도달 시점,
      이동평균 최종 높이).
  \end{itemize}
  \vskip0.5em
  \begin{block}{표준: 여러 시도의 평균 곡선}
    한 시드만 보고 ``안 된다/잘 된다''로 결론 =
    한 주가의 주가만 보고 ``나쁜 회사/좋은 회사''를 단정하는 것.
    \textbf{최소 3$\sim$5개 시드}로 \textbf{평균 $\pm$ 표준편차}를 보고,
    (1) ``이동평균이 \textbf{일관되게} 오르는가''(방향)와
    (2) ``시드 간 편차가 어느 정도인가''(안정성)를 \textbf{둘 다} 평가.
  \end{block}
  {\footnotesize Week 8.1 팀 프로젝트 FAQ의 ``한 실험으로 결론 내리지
  않는다''의 강화학습 버전.}
\end{frame}
